% Adapted from ne630_problems/lesson_16/statement.tex at commit e1fba66.
% Problem 1 of the source Lesson 16 assignment is used for Homework 15.

\noindent{\bf Problem 1 [FNRP 4.1].}

A reactor is to be built with fuel rods of diameter $D=1.2~\mathrm{cm}$ and
a liquid moderator with a 2:1 moderator-to-fuel volume ratio. What is the
distance $P$ between nearest fuel centerlines (a) for a square lattice and
(b) for a hexagonal lattice? State how cladding volume is treated.

\begin{adaptedsolution}
Neglect cladding volume (equivalently, assign it consistently to either fuel
or moderator). Per unit axial length,
\[
V_f=\frac{\pi D^2}{4},\qquad V_m=2V_f,
\]
so the unit-cell area is $V_f+V_m=3V_f$.

For a square lattice, the cell area is $P^2$:
\[
P=\sqrt{3V_f}=D\sqrt{\frac{3\pi}{4}}
=\boxed{1.84~\mathrm{cm}}.
\]
For a triangular lattice (a hexagonal Wigner--Seitz cell), the area per rod
is $\sqrt{3}P^2/2$:
\[
\frac{\sqrt{3}}{2}P^2=3V_f,
\qquad
P=D\sqrt{\frac{\sqrt{3}\pi}{2}}
=\boxed{1.98~\mathrm{cm}}.
\]
Both pitches exceed the rod diameter, as required.
\end{adaptedsolution}
